% Draft Introduction (rev. 2) --- gradual opening; pedology lower & cooler; no
% "paradigm shift"; hook = "emergent properties of SOC decomposition"; edaphic
% corroboration kept. Staging file for voice approval (2026-06-29).
% Reuses existing citation keys; [CITE: ...] marks gaps to fill.

\section{Introduction}
\label{sec:intro}

\paragraph{Soil-carbon turnover and the climate feedback.}
Soils store roughly twice as much carbon as the atmosphere, and the rate at which
that carbon cycles --- not the size of the stock alone --- sets much of the land
carbon--climate feedback. The relevant quantity is the mean transit time $\tau$,
the average time a carbon atom spends between entering and leaving the soil
\citep{Sierra2017,Sierra2018}. Because the soil stock is so large, even modest
shifts in $\tau$ across large areas reshape the global carbon balance: what
controls $\tau$, and how it varies from place to place, bears directly on whether
soils continue to buffer anthropogenic emissions or begin to amplify them, and is
thus a first-order term in the trajectory of atmospheric CO$_2$.

\paragraph{The classical paradigm: chemistry and constants.}
In the models that have represented this process since the 1970s--1990s
% [CITE: Jenkinson; Parton/CENTURY; RothC lineage]
, decomposition is first-order: the loss rate from each pool is proportional to
its mass, scaled by functions of temperature and moisture. The rate constant $k$
--- and thus $\tau = 1/k$ --- is treated as an intrinsic property of the
substrate, spatially invariant once the environmental scalars are applied. This
descends from a view in which substrate chemistry sets decomposability, and Earth
System Models (ESMs) still largely inherit the architecture
% [CITE: Todd-Brown 2013; Wieder 2013]
.

\paragraph{Persistence as an ecosystem property.}
This view has been progressively revised. \citet{Schmidt2011} and subsequent work
recast SOC persistence as an ecosystem property rather than a molecular one,
emerging from interactions among organic matter, mineral surfaces, microbial
communities, and soil structure: organo-mineral associations stabilise carbon
independently of its chemistry
% [CITE: Kleber 2007/2021; Lehmann \& Kleber 2015]
; microbial carbon-use efficiency varies with community composition
% [CITE: Kallenbach 2016; Cotrufo MEMS]
; mycorrhizal pathways shape the quantity and fate of belowground carbon
% [CITE: Averill 2014]
. Decomposition kinetics, in this light, are less fixed constants than the local
outcome of interacting ecological, edaphic, and biological processes.

\paragraph{A zonality of soil-carbon turnover.}
If the emergent properties of SOC decomposition arise locally in this way, they
should carry spatial structure. Here pedology offers more than an analogy: soils
are themselves organised zonally --- the soil-forming factors vary in coherent
geographic patterns, and soil types with them --- and an emergent property of the
soil system has reason to share part of that organisation. Under this hypothesis
$\tau$ is not a universal constant scaled by climate but carries a residual spatial
signal reflecting edaphic, biological, and land-use context --- structure that
climate scalars alone do not capture.

\paragraph{Climate and soil as coupled controls.}
Climate is itself a primary soil-forming factor, so edaphic properties and climate
are not independent: they carry overlapping signals about turnover. Where soil
properties predict $\tau$ about as well as climate does, and the two share much of
their explanatory power, that shared signal is an expression of the zonal coupling
that pedology describes. Read this way, the zonal perspective is not a competitor
to the climate-driven account but an alternative framing of it --- one that treats
climate as one soil-forming factor among several rather than the master variable,
and reads soil-carbon turnover as a property of the whole soil system.

\paragraph{Reading the structure rather than imposing it.}
One could try to represent these interactions mechanistically, but any such model
imposes a chosen structure on a system whose processes are incompletely known, and
risks encoding assumptions as results. We take the complementary route: rather
than prescribe how non-climate context acts, we ask what spatial structure the
observations already carry once climate is accounted for, and let that structure
--- where it is real, where it is organised, and how it behaves --- constrain
interpretation. The comparison with ESMs is the closing question, not the
foundation.

\paragraph{This study.}
We compile a global set of mineral-soil profiles, estimate steady-state transit
times, and interrogate the variance that remains once climate is accounted for
under spatial cross-validation. We ask \emph{how much} of it belongs to which
process class --- climate, edaphic, biological, land use --- attributed in an
order-independent way; whether it is \emph{organised} in space or merely noise;
\emph{where} non-climate context lengthens or shortens turnover, mapped at a scale
the data can support; and \emph{how} each driver acts. Finally we ask whether the
ESMs used for projection represent the resulting spatial structure or distort it.
The beyond-climate signal, though modest in size, \emph{has} a coherent and
ecologically interpretable structure --- edaphic-led, with soil biology as a
distinct second axis --- that ESMs amplify too steeply toward the poles and
disagree on among themselves by an order of magnitude.
